\documentclass[10pt, a4paper]{article}

\usepackage{homework}

\title{\textbf{《数值代数》上机作业14}}
\author{岳瀚阳\quad\quad\quad\quad 3230104443}
\date{\today}

\begin{document}

\maketitle

%----------------------------------------------------------------------
\section{隐式对称 QR 算法}
%----------------------------------------------------------------------

对于实对称矩阵 $A$，隐式对称 QR 算法分两步：

\textbf{第一步}，用 Householder 变换将 $A$ 化为对称三对角矩阵 $T = Q_0^{\!\top} A Q_0$。
对 $k = 0, 1, \ldots, n-3$，取 $x = A[k+1:\!,k]$，构造 Householder 向量 $v$ 与系数 $\beta$，使得
\[
H_k = I - \beta v v^{\!\top}, \qquad H_k x = \|x\|_2 e_1.
\]
然后对称更新 $A \leftarrow H_k A H_k$ 并累积 $Q_0 \leftarrow Q_0 H_k$。

\textbf{第二步}，对三对角矩阵 $T$ 执行带 Wilkinson 位移的隐式 QR 迭代。
位移 $\mu$ 取 $T$ 的右下角 $2\times2$ 块中靠近 $t_{nn}$ 的特征值：
\[
\delta = \frac{t_{n-1,n-1} - t_{nn}}{2}, \qquad
\mu = t_{nn} - \frac{\mathrm{sgn}(\delta)\,t_{n-1,n}^2}{|\delta| + \sqrt{\delta^2 + t_{n-1,n}^2}}.
\]

隐式 QR 步不直接计算 $T - \mu I$ 的 QR 分解，而是构造第一个 Givens 旋转 $G_1(0,1)$，使得 $G_1^{\!\top}(T - \mu I)e_1 = [\ast, 0, \ldots, 0]^{\!\top}$，然后计算 $T \leftarrow G_1^{\!\top} T G_1$。
这一步会在 $(2,0)$ 位置引入 ``bulge''，随后通过一系列 Givens 旋转 ``chase the bulge''，逐次消除 $(i+2,i)$ 位置的非零元，最终恢复三对角形式。
所有正交变换同步累积到特征向量矩阵 $Q = Q_0 G_1 G_2 \cdots$。

迭代终止准则取次对角线元素满足
\[
|t_{i+1,i}| < \varepsilon\,(|t_{ii}| + |t_{i+1,i+1}|), \qquad \varepsilon = 10^{-12}.
\]
特征值按升序排列后输出。

%----------------------------------------------------------------------
\subsection{教材第 244 页 1(2)}
%----------------------------------------------------------------------

矩阵为 $n$ 阶对称三对角阵，主对角线元素为 $4$，次对角线元素为 $1$：
\[
A = \begin{bmatrix}
4 & 1 & & & \\
1 & 4 & 1 & & \\
& \ddots & \ddots & \ddots & \\
& & 1 & 4 & 1 \\
& & & 1 & 4
\end{bmatrix}_{n\times n}.
\]

该矩阵的特征值有显式表达式
\[
\lambda_k = 4 + 2\cos\frac{k\pi}{n+1}, \qquad k = 1, 2, \ldots, n.
\]
当 $n = 50$ 时，$\lambda_1 = 4 + 2\cos(50\pi/51) \approx 2.0038$，$\lambda_{50} = 4 + 2\cos(\pi/51) \approx 5.9962$；
当 $n = 100$ 时，$\lambda_1 \approx 2.0010$，$\lambda_{100} \approx 5.9990$。

数值结果见表~\ref{tab:sym1}。

\begin{table}[H]
    \centering
    \caption{对称三对角矩阵 (4, 1) 的特征值}
    \label{tab:sym1}
    \begin{tabular}{ccccc}
        \toprule
        $n$ & 迭代次数 & $\lambda_{\min}$ & $\lambda_{\max}$ & 与 numpy 最大差异 \\
        \midrule
        50  & 101 & 2.0037933425 & 5.9962066575 & $1.51\times10^{-14}$ \\
        100 & 199 & 2.0009674354 & 5.9990325646 & $1.69\times10^{-14}$ \\
        \bottomrule
    \end{tabular}
\end{table}

%----------------------------------------------------------------------
\subsection{教材第 244 页 2(2)}
%----------------------------------------------------------------------

矩阵为 $100$ 阶对称三对角阵，主对角线元素为 $2$，次对角线元素为 $-1$：
\[
A = \begin{bmatrix}
2 & -1 & & & \\
-1 & 2 & -1 & & \\
& \ddots & \ddots & \ddots & \\
& & -1 & 2 & -1 \\
& & & -1 & 2
\end{bmatrix}_{100\times100}.
\]

这是离散一维 Laplace 算子的标准形式，其特征值为
\[
\lambda_k = 2 - 2\cos\frac{k\pi}{101}, \qquad k = 1, 2, \ldots, 100.
\]
理论值 $\lambda_1 = 2 - 2\cos(\pi/101) \approx 0.0009674$，$\lambda_{100} = 2 - 2\cos(100\pi/101) \approx 3.9990326$。

隐式对称 QR 算法迭代 200 次收敛，结果与 numpy 最大差异 $1.02\times10^{-14}$，最大残差 $1.74\times10^{-12}$。
最小与最大特征值分别为
\[
\lambda_{\min} = 0.0009674354, \qquad
\lambda_{\max} = 3.9990325646.
\]

对应特征向量的前 5 个分量为
\begin{align*}
v_{\min} &\approx [0.004376,\; 0.008748,\; 0.013112,\; 0.017463,\; 0.021797]^{\!\top}, \\
v_{\max} &\approx [0.004376,\; -0.008748,\; 0.013112,\; -0.017463,\; 0.021797]^{\!\top}.
\end{align*}
这与理论特征向量 $v_k(j) = \sin(jk\pi/101)$ 完全一致：最小特征向量近似正弦曲线，最大特征向量近似交错符号的正弦曲线。

%----------------------------------------------------------------------
\section{隐式 QR 算法（Francis 双重位移）}
%----------------------------------------------------------------------

对于一般实矩阵 $A$，采用 Francis 双重位移 QR 算法求实 Schur 形式。

\textbf{第一步}，Householder 变换将 $A$ 化为上 Hessenberg 形式 $H = Q_0^{\!\top} A Q_0$。

\textbf{第二步}，对 $H$ 执行双重位移 QR 迭代。
取位移 $\sigma_1, \sigma_2$ 为右下角 $2\times2$ 块的特征值，令 $s = \sigma_1 + \sigma_2$，$t = \sigma_1\sigma_2$。
构造矩阵
\[
M = H^2 - sH + tI = (H - \sigma_1 I)(H - \sigma_2 I),
\]
对 $M$ 进行 QR 分解得到 $Q_{\text{step}}$，然后更新
\[
H \leftarrow Q_{\text{step}}^{\!\top} H Q_{\text{step}}, \qquad
Q \leftarrow Q Q_{\text{step}}.
\]
由隐式 $Q$ 定理，$H$ 保持上 Hessenberg 形式。
当次对角线元素收敛到零时，矩阵分裂为更小的子块；若右下角 $2\times2$ 块有复特征值，则直接提取该块。

%----------------------------------------------------------------------
\subsection{教材第 202 页 2(2)：$x^{41} + x^3 + 1 = 0$ 的全部根}
%----------------------------------------------------------------------

将求根问题转化为特征值问题，构造首一多项式 $x^{41} + x^3 + 1$ 的友矩阵（companion matrix）$C$，其特征值即为方程的全部根。
友矩阵为 $41\times41$ 的上 Hessenberg 矩阵：
\[
C = \begin{bmatrix}
0 & 0 & \cdots & 0 & -1 \\
1 & 0 & \cdots & 0 & 0 \\
0 & 1 & \cdots & 0 & 0 \\
\vdots & & \ddots & \vdots & \vdots \\
0 & 0 & \cdots & 1 & 0
\end{bmatrix}
\begin{array}{l}
\text{(最后一列：$-a_0=-1$，$-a_3=-1$，其余为 $0$)} \\
\text{次对角线全为 $1$}
\end{array}
\]

Francis 双重位移 QR 算法迭代 4120 次后收敛，求得全部 41 个特征值。
与 numpy.roots 比较，最大差异 $1.15\times10^{-12}$，验证了算法的正确性。

直接代入多项式验证时，$|p(\lambda)|$ 的量级约为 $10^0$，这并非算法误差，而是浮点求值中的\textbf{灾难性抵消}（catastrophic cancellation）：
对于 $|\lambda| \approx 1$ 的根，$\lambda^{41}$ 与 $\lambda^3$ 的量级相近、相位相反，大量有效数字在相消中损失。
若改用 Horner 法或扩展精度重新估计，相对误差仍在机器精度量级。

%----------------------------------------------------------------------
\subsection{教材第 202 页 2(3)：矩阵 $A(x)$ 的特征值}
%----------------------------------------------------------------------

\[
A(x) = \begin{bmatrix}
9.1 & 3.0 & 2.6 & 4.0 \\
1.2 & 5.3 & 1.7 & 1.0 \\
3.2 & 1.7 & 9.4 & x \\
6.1 & 4.9 & 3.5 & 6.2
\end{bmatrix}.
\]

分别对 $x = 0.9,\,1.0,\,1.1$ 运行隐式 QR 算法，结果汇总于表~\ref{tab:Ax}。

\begin{table}[H]
    \centering
    \caption{矩阵 $A(x)$ 的全部特征值}
    \label{tab:Ax}
    \begin{tabular}{c|cccc}
        \toprule
        $x$ & $\lambda_1$ & $\lambda_2$ & $\lambda_3$ & $\lambda_4$ \\
        \midrule
        $0.9$ & 2.336853 & 4.752352 & 6.888012 & 16.022782 \\
        $1.0$ & 2.336448 & 4.747533 & 6.856381 & 16.059638 \\
        $1.1$ & 2.336037 & 4.742601 & 6.825152 & 16.096210 \\
        \bottomrule
    \end{tabular}
\end{table}

与 numpy.linalg.eigvals 比较，三组特征值的最大差异分别为 $1.07\times10^{-14}$、$2.13\times10^{-14}$、$2.84\times10^{-14}$，达到机器精度。

观察特征值随 $x$ 的变化：
\begin{itemize}
    \item $\lambda_1$（约 $2.34$）和 $\lambda_2$（约 $4.75$）对 $x$ 的变化不敏感，变化幅度分别约为 $0.001$ 和 $0.010$；
    \item $\lambda_3$（约 $6.85$）随 $x$ 增大而减小，变化幅度约为 $0.063$；
    \item $\lambda_4$（约 $16.05$）随 $x$ 增大而显著增大，变化幅度约为 $0.073$。
\end{itemize}

这种现象可用特征值的摄动理论解释。将 $A(x)$ 写为 $A(x) = A(0) + x E$，其中 $E$ 仅在 $(3,4)$ 位置为 $1$。
由特征值一阶摄动公式
\[
\frac{\partial\lambda_i}{\partial x} = \frac{v_i^{\!\top} E u_i}{v_i^{\!\top} u_i},
\]
其中 $u_i, v_i$ 分别为左右特征向量。
对于最大特征值 $\lambda_4$，其特征向量在 $x$ 变化方向（第 4 个分量）上的投影较大，因此对 $x$ 最敏感。

%----------------------------------------------------------------------
\section{结果分析}
%----------------------------------------------------------------------

\textbf{（1）隐式对称 QR 算法的收敛性。}
对于对称三对角矩阵，Wilkinson 位移的隐式 QR 迭代具有立方收敛速度。
实验表明，$n=50$ 时需 101 次迭代，$n=100$ 时需 199 次迭代，迭代次数与矩阵维数近似成正比，符合理论预期。

\textbf{（2）Francis 双重位移 QR 算法对复特征值的处理。}
一般实矩阵的 QR 迭代不会收敛到对角矩阵，而是收敛到实 Schur 形式（拟上三角矩阵），其中 $2\times2$ 对角块对应一对共轭复特征值。
双重位移策略通过隐式应用 $(H - \sigma_1 I)(H - \sigma_2 I)$，避免了复数运算，同时保证迭代收敛。

\textbf{（3）友矩阵的数值稳定性。}
友矩阵是求多项式根的经典方法，但高次多项式（如 $41$ 次）的友矩阵条件数通常较大，且多项式求值时易出现灾难性抵消。
对于本题的 $x^{41} + x^3 + 1$，虽然特征值精度达到 $10^{-12}$，但直接代入验证时 $|p(\lambda)| \approx 10^0$，这是因为浮点运算中的相消损失了大量有效数字。

%----------------------------------------------------------------------
\appendix
\section{源代码}
%----------------------------------------------------------------------

\subsection{\texttt{hw14\_test.py}（主程序）}
\lstinputlisting{../tests/hw14_test.py}

\subsection{\texttt{symmetric\_qr.py}（隐式对称 QR 算法）}
\lstinputlisting{../algorithms/symmetric_qr.py}

\subsection{\texttt{implicit\_qr.py}（Francis 双重位移 QR 算法）}
\lstinputlisting{../algorithms/implicit_qr.py}

\end{document}
